\section{Related Work}
\label{sec:related}

\textbf{Heuristic Computer Vision in Agriculture:} The earliest attempts to automate crop monitoring utilized classical computer vision paradigms, which depended on manually engineered features like color thresholding and edge maps. These systems were rigid and struggled to generalize across diverse field conditions. Variations in weather, time of day, and crop overlap caused these primitive algorithms to break down, resulting in massive false-positive rates and unreliable yield estimations.

\textbf{The Shift to Convolutional Networks:} The deep learning revolution brought Convolutional Neural Networks (CNNs) to the forefront of object detection. Architectures utilizing a two-stage approach, such as R-CNN \cite{girshick2014rich} and its successors (Fast and Faster R-CNN), significantly advanced precision by separating region proposal generation from the final classification step. While these networks established a high ceiling for accuracy, their sequential nature incurred substantial latency. This computational burden renders two-stage detectors highly impractical for deployment on the embedded edge devices typically found on agricultural drones.

\textbf{Real-Time Inference and Attention Mechanisms:} To bypass the latency bottleneck, single-stage detectors fundamentally altered the detection paradigm. The YOLO (You Only Look Once) framework \cite{redmon2016you} unified region extraction and classification into a single network pass, achieving breakthrough inference speeds. Although early iterations struggled to resolve small, clustered objects (like wheat heads), modern enhancements like multi-scale feature pyramids have resolved these spatial limitations. Conversely, while Vision Transformers \cite{vaswani2017attention} and attention-based models like DETR \cite{carion2020end} have pushed the boundaries of theoretical accuracy, their intense computational demands make them unsuitable for real-time agricultural applications.

\textbf{The CropSense Methodology:} In developing our framework, we prioritized the optimal balance between high-fidelity detection and rapid inference speeds. Heavyweight transformers and two-stage networks are too slow for real-time UAV integration. Consequently, we anchored our approach on the highly optimized YOLOv8 architecture, tuning its hyperparameters specifically to resolve the dense, minute geometries characteristic of the Global Wheat dataset, ensuring real-time performance without sacrificing analytical precision.